\documentclass{article}
\usepackage[utf8]{inputenc}
\usepackage[T1]{fontenc}
\usepackage{geometry}
\usepackage{graphicx}
\usepackage{float}
\usepackage{hyperref}
\usepackage[english,indonesian]{babel}

\title{Optimasi Pemrosesan Template LaTeX Menggunakan Pendekatan Berbasis AI}
\author{Academic Support Agent}
\date{\today}

\begin{document}
\maketitle

\begin{abstract}

\end{abstract}

\chapter{PENDAHULUAN}


\chapter{TINJAUAN PUSTAKA}


% Original content (commented out):
% %%%%%%%% DO NOT MODIFY %%%%%%%
\documentclass[11pt]{article}
\usepackage{tls}
%%%%%%%%%%%%%%%%%%%%%%%%%%%%%

% Any custom packages you wanna use go below this line
% \usepackage{package-name}

% Edit the title, author names and your current email addresses
\title{\liningfont Template for submitting to TLS Proceedings}
\author{\liningfont Author 1 \quad Author 2 \\ University 1 \quad University 2 \\ \texttt{email@univ.edu} \quad \texttt{email@univ.edu}}
\date{}           % Leave date empty

\begin{document}

% DON'T COMMENT/REMOVE THE FOLLOWING 3 LINES
\maketitle
\thispagestyle{empty}
\pagestyle{empty}

% START WRITING YOUR CONTENT BELOW THIS LINE

\begin{abstract}
\noindent This document both describes and exemplifies the TLS proceedings format. You must use this \LaTeX{}  source as a template for your submission. The sections below specify requirements for the paper and supply examples for how to typeset common linguistic phenomena in \LaTeX{}, as well as citation best practises. This document was updated in Spring 2023.
\end{abstract}

\section{Introduction}

The following rules apply to all submitted papers:

\begin{itemize}
    \item They must be written in English.
    \item They must contain the name(s) of the author(s).
    \item They cannot exceed 20 pages, including references, data, and appendices.
\end{itemize}

You must submit both the final PDF as well as the original \LaTeX{} source files (\texttt{.tex} and \texttt{.bib} files, as well as all your figures). This paper template can be found on the conference website, as well as on Overleaf\footnote{Search for `TLS' at \url{https://www.overleaf.com/latex/templates}}. \textbf{We recommend you use the template on Overleaf} --- Overleaf is fast, free, and easy to use, and takes away the complexity of LaTeX installation. You can make a copy of this template on Overleaf when signed in and edit to make your final submission. If there are special questions or wishes regarding paper preparation and submission for TLS, correspondence should be addressed via email to \url{tls.conference@gmail.com}; Please include `TLS [YEAR] Proceedings' in the subject line of the email.

\paragraph{Overleaf settings} You should be able to compile your \TeX{} source file with the default settings. Ensure that `Compiler' in the Overleaf menu (opened by clicking on the Menu button in the top left of your browser window) is set to \textbf{LuaLaTeX}.

\section{Page Layout and Style}

The page layout should conform to the following rules. The easiest way to meet these requirements is to use the supplied \LaTeX{} template and check details against this example file. Please follow these rules as carefully as possible and contact the editors for any clarifications.

\subsection{Basic Layout}

Please adhere to the following basic layout parameters:

\begin{itemize}
    \item Page format should be A4.
    \item Left and right margins are 1.5''.
    \item Top and bottom margins are 1''.
    \item The header and footer of each page should be empty; page numbers and author/title headers will be added later by the TLS editors.
\end{itemize}

The \LaTeX{} template has the margins and page format setup correctly by default --- don't add headers or footers to the template when modifying it to write your paper.

\subsubsection{Headings}

Section headings (including sub-headings and sub-sub-headings) are left justified in bold-face with the first letter capitalized and the rest of the heading in lower case. No more than three levels of headings should be used. In \LaTeX{}, you can start a new section heading with \verb|\section{section-heading}| command, a new subsection with \verb|\subsection{subsection-heading}|, and a new sub-subsection with \verb|\subsubsection{subsubsection-heading}|.

\subsection{Font Faces}

Charter is used for the main body text using the \texttt{XCharter} package. Since we use the \texttt{tipa} package for typesetting IPA characters, IPA is set using the Charis SIL font. Both fonts are from the same family, offer a broad coverage of most characters in Unicode, and should allow you to typeset characters in most languages.

\paragraph{Multilingual typesetting} Please refer to the Overleaf documentation\footnote{\url{https://www.overleaf.com/learn}} on how to use multiple fonts in your document for languages other than English. Charter offers a broad coverage of languages and characters, but if you find some characters not appearing properly, please refer to the Overleaf documentation for a list of fonts supported natively by Overleaf in that language. If you're still having trouble, contact the TLS editors.

\subsection{Figures}

All figures should be centered on the page. Figure captions should follow each figure and have the format given in Figure~\ref{fig:label}.
Figures should preferably be line drawings. If they contain gray shades, it should be checked that they print well on a high-quality non-color laser printer. Color figures should not be used.

Adjust the width of the image using \verb|\columnwidth| parameter. For instance, to set the image to be half the width of the page, use \verb|0.5\columnwidth|. You can also adjust the placement of the figure with the float parameter when starting a new figure: \verb|t| says to place the figure at the top of a page, while \verb|H| will place it in the position where you insert the above markup into your \LaTeX{} source file.

To refer to the figure in text, assign a label to the figure in \LaTeX{}, as described in the example with \verb|\label{fig:label}|. Then you can refer to the figure number with \verb|\ref{fig:label}| in your text. You can also use this same labelling procedure to refer to tables and examples.

\begin{figure}[t]
    \centering
    \includegraphics[width=0.5\columnwidth]{figures/vowel-chart.png}
    \caption{IPA vowel chart}
    \label{fig:label}
\end{figure}  


\subsection{Tables}

An example of a table is shown in Table 1. Somewhat different styles are allowed according to the type and purpose of the table. Color should not be used, but gray shading is allowed.

\input{tables/table.tex}

\subsection{Examples}

We recommend using the gb4e package. It’s not perfect, but it allows for sub-examples (which are not supported by the equation environment), and interlinear glosses:

\begin{exe}
    \ex 
    \begin{xlist} 
        \ex
        \gll Los niños le molestan \\ the.{\scshape pl} children {\scshape dat} bother\\
        \glt `Children irritate him'
        \ex
        \gll aux chevaux\\ to.{\scshape art.pl} horse.{\scshape pl}\\
        \glt `to the horses'
    \end{xlist} 
    
    \ex \gll My s Marko poexa-l-i avtobus-om v Peredelkino. \\ 
    {\scshape 1pl} {\scshape com} Marko go-{\scshape pst-pl} bus-{\scshape ins} All Peredelkino.\\
    \glt `Marko and I went to Perdelkino by bus.'
\end{exe}

To keep your glosses lined up properly we recommend the \texttt{exe} and \texttt{gll} macros from the \texttt{gb4e} package.

\subsubsection{Equations}

On \LaTeX{}, do not use the \texttt{equation} macro, as it uses a different counter that does not integrate with the \texttt{gb4e} counter. To typeset a formula in math mode, use the \$\textellipsis \$ delimiters. To display the equation center aligned in math mode, use the \verb|center| macro:

\begin{center}
    $t_0 = \frac{1}{f_0}$
\end{center}


\subsection{OT Tableaux}

To generate Optimality Theory (OT) tableaux, please refer to the documentation for the \texttt{ot-tableaux} package at \url{https://ctan.org/pkg/ot-tableau}. Remember to load the package at the beginning of the \LaTeX{} file if you plan on using it.


\subsection{Syntactic Trees}

We recommend using the \texttt{tikz-qtree} package for generating syntactic trees. For example, Figure~\ref{fig:tree} is a simple tree generated using \texttt{tikz-qtree} that illustrates movement.

\begin{figure}[t]
        \centering
        \footnotesize
        \begin{tikzpicture}
            \tikzset{every tree node/.style={align=center,anchor=north}}
            \Tree[.TP [.DP \edge[roof]; {The gardener} ] [.T$'$ [.T T\\will$_{[\textrm{+fin:fut, \checkmark T}]}$ \node (a) {V\\have$_{i,[\textrm{\checkmark \textit{u}T}]}$}; ] [.VP$'$ \phantom{} [.V$'$ [.AdvP \edge[roof]; {really} ] [.V$'$ \node (b) {V\\$t_i$}; [.VP$'$ \edge[roof]; {frightened those grackles} ] ] ] ] ] ]
        \draw [->] (b) [in=-90,out=-90,looseness=2]  to (a);
        \end{tikzpicture}
        \caption{A syntactic tree with movement.}
        \label{fig:tree}
\end{figure}


\subsection{References}

We use the LSA style of \texttt{name (year)} or \texttt{(name year)}. The following citation commands can be used in \LaTeX{} to get different citation styles, depending on your need:

\begin{itemize}
    \item \verb|\parencite{England1991}| produces a parenthetical citation, which is displayed as: \parencite{England1991}
    \item \verb|\textcite{Beavers2009}| produces a text citation, which is displayed as: \textcite{Beavers2009}
    \item \verb|\footcite{Erk2008}| produces a footnote citation, which is displayed as a footnote~\footcite{Erk2008}.
\end{itemize}

Try to include links to the papers that you cite in your submission by adding the \verb|howpublished={\url{link-to-paper}}| line to your \texttt{bib} entry.

\subsection{Hyperlinks}

URLs, email addresses, etc. should be displayed in full and formatted using the \verb|\url{}| macro. For instance, \verb|\url{http://tls.ling.utexas.edu}| in \LaTeX{} is displayed as \url{http://tls.ling.utexas.edu}. This will typeset the URL in monospaced font with a blue color.

Please don't use inline links using \verb|\href| macro from the \verb|hyperref| --- they will not show up in the final version of your document in the published proceedings.

\subsection{Footnotes}

Footnotes can be appended using the \verb|\footnote{...}| macro.\footnote{Remember that these are \emph{footnotes}, not places to go off on multi-paragraph rants.}

\subsection{Bibliography}

The \verb|\printbibliography| macro at the end of the \LaTeX{} source document automatically writes your bibliography with proper formatting at the end of your PDF submission.

If you have any questions with writing in \LaTeX{}, refer to the Overleaf documentation, the source for this template, or email the TLS editors.

% Print Bibliography
\printbibliography

\end{document}


\end{document}